%% kapa: exam=1 round=24 year=2013 period=2 form=A subject=appraisal_law subject_ko=부동산관계법규 seq=40 number=80
%% source: Q-Net 감정평가사 기출문제 (cst00309) artlSeq=5073802 file=extracted/1008/5073802/2013년 제24회 감정평가사 1차시험 기출문제/2교시(A형).hwp
%% conversion: kordoc 4.24.51 -> tools/gen1.py (auto); review=proofread-sonnet; flags=none
\begin{question}{80}
부동산등기법상 등기를 할 때 부기로 하여야 하는 것을 모두 고른 것은?
\begin{framedtext}
ㄱ. 환매권실행등기 \par
ㄴ. 권리소멸약정등기 \par
ㄷ. 말소회복등기 \par
ㄹ. 소유권 외의 권리의 이전등기 \par
ㅁ. 부속건물의 신축등기 \par
ㅂ. 등기명의인표시의 변경등기
\end{framedtext}
\begin{choices}
  \item ㄱ, ㄴ, ㄷ
  \item ㄱ, ㄷ, ㅁ
  \item ㄴ, ㄷ, ㄹ
  \item ㄴ, ㄹ, ㅂ
  \item ㄹ, ㅁ, ㅂ
\end{choices}
\end{question}
\answer{4}
